% =====================================================================
%  DECISION RECORD — what was merged into aluminium_report_downstream.tex
%  in response to the reviewer comments and the employment-per-tonne note,
%  what was deliberately left out, and what remains outstanding.
% =====================================================================
\documentclass[11pt, a4paper]{article}
\usepackage[a4paper, top=2.2cm, bottom=2.2cm, left=2.2cm, right=2.2cm]{geometry}
\usepackage{fontspec}
\IfFontExistsTF{DejaVu Sans}
  {\setmainfont{DejaVu Sans}}
  {\IfFontExistsTF{Verdana}{\setmainfont{Verdana}}{\setmainfont{Arial}}}
\usepackage{booktabs}
\usepackage{amsmath}
\usepackage{enumitem}
\usepackage{titlesec}
\usepackage{float}
\emergencystretch=2.5em
\hbadness=10000
\titleformat{\section}{\large\bfseries}{\thesection}{0.7em}{}

\title{\textbf{Decision Record: Reviewer Comments and the Per-Tonne Note}\\
\large What was merged, what was not, and why}
\author{}
\date{}

\begin{document}
\maketitle
\vspace{-1.0cm}

\section{Merged into the report}

Everything below is now in \texttt{aluminium\_report\_downstream.tex}. Each item is
either arithmetic derived from data already in the report, or a claim with a
citation that has been independently verified.

\begin{table}[H]
\centering
\small
\begin{tabular}{p{4.3cm}p{2.2cm}p{8.4cm}}
\toprule
\textbf{Comment} & \textbf{Where} & \textbf{What was added} \\
\midrule
Is causation established? & §4.1 & A paragraph stating that the evidence is
association, not an identified causal effect: no counterfactual, three years of ERP
against twelve of ASI, construction cycle uncontrolled. Two overclaiming sentences
in §4.6 and §4.7 were rewritten. \\
\addlinespace
Effect of a rise in metal cost & §4.7, App.~B.4 &
$\Delta V/V = -[m\theta/(1-m)]\hat p$. With $m = 0.831$ the amplification factor is
\textbf{4.92}; at $\theta = 60\%$ a 1\% rise in metal price removes $\approx$3\% of
downstream value added and the 7.5\% duty over a fifth of it. Upstream the figure
is nil. \\
\addlinespace
Employment per tonne & §4.3, App.~B.4 &
\textbf{41 persons per thousand tonnes upstream against 210--262 downstream}
(5.1--6.4), with the exact decomposition
$3.14 \times 1.81 = 5.7$ linking it to the per-rupee ratio. \\
\addlinespace
Export mix (from the note) & §3.1 &
India's exports are \textbf{61.4\% primary metal}; China's are \textbf{2.8\%}.
India imported US\$4.1bn of finished aluminium goods in FY2025--26, a quarter at
low or zero duty. Sourced to GTRI, added to the bibliography. \\
\addlinespace
Actual tax or a proxy? & §4.7 &
An explicit statement that corporate tax is an \emph{estimate} of the tax base, not
receipts; axis and table labels renamed accordingly. \\
\addlinespace
``Control'' is wrong for the UK & §3.3 &
Rewritten as a non-preferential (MFN) benchmark holding product basket and estimator
constant, with the explicit note that no further causal weight is placed on it. \\
\addlinespace
Price gap not shown to be tariff-driven & §3.1 &
Softened: the duty is ``one mechanism through which such a spread can be sustained'',
not its demonstrated cause, with the missing decomposition named. \\
\bottomrule
\end{tabular}
\end{table}

\section{Deliberately not merged}

\begin{enumerate}[leftmargin=1.5em, itemsep=4pt]

\item \textbf{The 35:1 employment ratio.} It compares the direct workforce of the
primary smelters alone (4.6 times smaller than the ASI's count for NIC~2420) with the
direct, indirect and unorganised workforce downstream (1.7 times larger than the ASI's
count for the fabrication classes). Those two mismatches multiplied by the
like-for-like ratio of 5.7 give $\approx$45, which is where the figure comes from. The
report uses 5.1--6.4 and explains the larger number in a footnote, so that a reader
who has seen it elsewhere is not left with an apparent contradiction.

\item \textbf{The stylised 2,000-tonne table.} Its parameters contradict the real
figures cited two paragraphs later in the same note (90 workers per thousand tonnes
upstream, against 9--10 in the same document). Two of its four summary percentages are
wrong: 150 against 180 workers is 83\%, not 67\%; \$70 against \$25 per tonne is 180\%
more, not 40\%. And its value-added-per-worker figures imply downstream is \emph{less}
labour-intensive per rupee---the opposite of what the ASI shows.

\item \textbf{``Immediate'' contract-worker adjustment.} Twelve annual observations
against three years of ERP estimates identify no lag structure, and the contract-share
series in fact moves gradually. The claim has been removed rather than defended; the
composition finding that replaces it needs no timing assumption.

\item \textbf{The industry's specific asks.} Drafted text exists but the asks and
their ordering are yours to set, not mine to infer from the report. Say the word and
it goes in.

\item \textbf{Stakeholder impact analysis.} A reasoned draft exists. It is
argumentation rather than evidence, and belongs in the Conclusion, which is outside
the section rebuilt so far --- flagging it rather than inserting it unilaterally.

\end{enumerate}

\section{Outstanding --- needs data}

\begin{enumerate}[leftmargin=1.5em, itemsep=4pt]
\item \textbf{Bilateral imports, HS 7604--7616 at eight digits, by partner, five
years} (DGCI\&S or UN~Comtrade). Would answer ``which Chinese products displace
Indian ones'', which the tariff files cannot --- they carry rates, not volumes. It
would also put the trade-remedy argument in §3.2 on evidence rather than assertion.
\emph{Highest value of the three.}

\item \textbf{Cost sheets from two or three representative fabricators}, giving metal
as a share of total input cost. This is $\theta$ in the amplification result, and the
only unmeasured parameter in the strongest quantitative claim in Section~4. It would
replace a range with a number.

\item \textbf{LME cash settlement, Indian ex-works prices, freight, insurance and
physical premia} over a stated window. Would let the price-gap claim be made properly
instead of softened.
\end{enumerate}

\end{document}
